\documentclass[12pt]{article}
\usepackage[a4paper, top=2.5cm, bottom=2.5cm, left=2cm, right=2cm]{geometry}

\usepackage{setspace}
\usepackage{graphicx}
\usepackage{booktabs}
\usepackage{hyperref}
\usepackage{amsmath}
%\usepackage[authordate,backend=biber,ibidtracker=false,uniquename=false,maxcitenames=3,dashed=false]{biblatex-chicago}

\usepackage{natbib}
\bibpunct{}{}{;}{a}{,}{,}

\usepackage{float}

%\addbibresource{Policing.bib}

\title{The Optics of Reform: Diversification and Trust in Policing}
\author{Eric Dobbie}
\date{\today}

\begin{document}
\maketitle
\doublespacing

\begin{abstract}
This study investigates how racial and gender diversity in policing shapes public trust, perceptions of fairness, and willingness to cooperate with law enforcement. While previous work has shown that racial and gender representation levels matter at extreme levels, there is lit Using a preregistered survey experiment, the study isolates the causal effects of representative cues on perceived legitimacy and behavioral intentions.
\end{abstract}

\newpage
\section{Introduction}

Representation is important for bureaucratic performance as the descriptive background of bureaucratic agent helps to shape their worldview and, as a consequence, their behavior (\cite{kennedy2017race, shoub2021female, ba2021role}). Police, in particular, are highly visible bureaucratic actors, and their behavior is a matter of public interest as police-civilian interactions are among the most common interactions between mass publics and the government. 

During a 2015 interview with BBC News, then Norfolk Police Department chief Michael Goldsmith described the importance his department needed to recruit Black police officers. ``If you don't have a diverse, inclusive workforce, you end up not making good decisions because you don't have everybody's perspective at the table.'' Norfolk PD was aiming to increase diversity in a police department with 16\% Black officers serving a population that was 43\% Black. (\cite{NorfolkPD_BBCRecruitment_2015})\footnote{Despite the announced effort to recruit more Black officers, according to the 2020 LEMAS survey, only 19\% of officers were identified as Black following Norfolk's recruiting effort  (\cite{LEMAS2020})}. Despite the evidence that representation effects outcomes, there is scarce work available on the effect of representation on trust in institutions.

Trust in police shapes perspectives and views of police legitimacy. As trust in police increases, improvements in police community relations can lead to better police work through information updating (\cite{mccall_resident_2019}). Many reform efforts are explicitly aimed at improving public trust in policing, mending strained relations between the community and police. These efforts include community oversight, technology, and training programs, all with the aim of improving policing outcomes and, especially in the case of community oversight, improving trust in policing. An alternative approach to improve police legitimacy is to change the \textit{who} the police are (\cite{meier2006gender, nicholson2017will}). By changing who makes up police departments through modified hiring practices police departments can change public perceptions.++

\section{Police Diversity and Trust}

Previous work on the effect of diversity in police departments have largely focused on the \textit{level} of diversity in police departments. \cite{stauffer2023police} utilize a conjoint experiment, varying both the level of non-white officers and women officers in respondents hypothetical police departments. The authors also vary department performance indicators, such as the felony solve rate, civilian complaints of police misconduct, rate of deadly force, rate of use of force, resources allocated to community policing, local budget allocated to force, and population served. They find that public trust was improved in a subset of populations, with improved perceptions of legitimacy among women as representation of women reaches the most intense treatment levels. Likewise, Black respondents view departments with a majority non-white officers are seen as more trustworthy and accountable.

Rather than focusing on the level of diversity in a police department, I focus on changes to police diversity. The theoretical foundations for focusing on changes rather than absolute levels follow from the behavior economics literature on reference dependence (\cite{kHoszegi2006model}). Rather than evaluating departments in absolute terms, individuals will evaluate using a reference point based on their expectations. Among typically excluded and marginalized populations, these expectations around police department composition and makeup are relatively fixed. Members of these marginalized group are also the least likely to trust the police (\cite{stauffer2023police}), so changes in the level of descriptive representation have the potential to increase their levels of institutional trust.


\subsection{Descriptive Representation}
In classic studies of descriptive representation, the focus has been on elected officials rather than bureaucratic actors (CITE SOME STUFF). This distinction makes sense as a starting point, as the mass public has the ability to  participate in the selection of these officials through elections. In contrast, bureaucrats are selected, in the case of police departments, through more standard recruiting and hiring practice. The public can only change these practices indirectly, through elections for city council or public pressure.

\subsection{Procedural Justice and Legitimacy}
Procedural justice, or the expectation of citizens that they will be treated fairly, is important for police legitimacy (\cite{hinds2007public, tyler2004enhancing}). Procedural justice can usually be broken down into four core components: (1) the expectation that citizens will participate before authorities make a decision, (2) perceived neutrality in decision making, (3) the authority demonstrated respect for the citizen during interactions, and (4) the authority conveys that they have good motives (\cite{goodman2010four,mazerolle2013procedural, tyler2002trust, tyler2011procedural}). \citet{mazerolle2013procedural} conducted a review of the literature, showing that, across 28 studies of procedural justice outcomes of satisfaction improved when police used tactics that align with the components of procedural justice. 

I argue that institutional level trust in police increase as descriptive representatives are added to the department because of expectations of procedural justice.

\subsection{Social Identity and Symbolic Representation}
This subsection draws on social psychology to explain why descriptive characteristics may serve as powerful cues for trustworthiness and fairness.

\subsection{Expectations}

Based on these theories, the following expectations can be derived. I expect that Black respondents will report higher levels of trust when the frame they are presented signals that the police department will hire more Black officers than in the control, general hiring case. I similarly expect that women will report higher levels of trust in departments that signal their intent to join the 30x30 program. 

\begin{quote}
    \textbf{Hypothesis 1:} Black respondents will report higher levels of trust for the racial diversity frame compared to the control baseline.\\
    \textbf{Hypothesis 2:} Women respondents will report higher levels of trust for the increase in female police officers frame compared to the control baseline.
\end{quote}



\section{Pilot Study}

I plan to run a pilot study to assess the uptake of treatment and calculate minimum detectable effects. In the pilot study, I plan to present respondents with the treatment vignettes to evaluate uptake of treatment and establish minimum detectable effects.

\subsection{Minimum Detectable Effects}

Based on the results of the pilot study, I plan to run a minimum detectable effect (MDE) analysis to ensure that the study can detect a meaningful effect for the given sample size. By identifying the minimum detectable effect, I will be better able to interpret potential null results given the minimum detectable effect. Using the MDE and the desired power of 0.8, I can calculate the necessary sample size for sufficiently powered analysis.


\section{Experimental Design}

I will utilize a between-subjects, block randomized experimental design with two treatment arms to test for public reactions to changes in police department hiring practices. Each participant will only observe one news article and evaluate that department. In the control condition, participants will be shown a news article describing a hiring drive that is aimed at increasing the number of officers employed by the department. There will be no emphasis on race or gender in the control condition. In the first treatment condition, the article will include a quote where the police chief expresses interest in increasing the number of Black officers during the recruiting drive to align the department more closely with the demographics of the community. This condition will not mention officer gender. In the final treatment arm, the chief will emphasize the goal of adding female officers to the force as a part of the 30x30 national campaign. This condition will not mention officer race. The remaining text, outside of the focus of the recruitment drive, will remain constant across the conditions to the extent practicable to maintain internal validity.

\subsection{Sampling and Recruitment}

I plan to field this survey through Prolific with a target sample size of 1,200 respondents. The study population is American adults. I plan to oversample Black respondents to make sure that conditional average treatment effect analysis for H1 is sufficiently powered.

\subsection{Randomization and Treatment Assignment}

I will utilize a block randomization approach. In block randomization, respondents are assigned to strata, or blocks, based on prognostic covariates, in this case race and gender, and assigned to treatment within blocks. This ensures a balanced sample across all blocking variables in the treatment and control conditions. Participants within each block will be assigned to each treatment arm with equal probability.

Respondents read a manipulated news article about a department’s hiring strategy, varying the focus of the recruiting drive. An example of the treatment vignettes is shown in Figure XX. All conditions can be found in Appendix XX.

In order to best reflect real life conditions, all treatment vignettes have been modified based on real articles describing hiring efforts of police departments. 

\subsection{Outcome Measures}

I use a set of three outcome measures. First, I ask respondents to rate their general level of trust in the police department to conduct day-to-day policing activities, such as traffic citations and beat work. Second, I ask respondents to rate their personal willingness to engage with the police department should they find themselves in physical danger. Third, and finally, I ask respondents to rate their likelihood of encouraging a friend of family member to apply to the department. Exact question wording is available in the survey instrument in Appendix XX.

These outcome measures each capture a different component of trust in policing. The first question aims to capture trust in the police to engage in routine public safety activity. 

The second question asks the respondent to put themselves in a situation where police intervention could be helpful and assesses their willingness to engage with this particular department. 

The third question aims to measure how comfortable respondents are with a friend or family member joining the department. 

\section{Empirical Strategy}

\subsection{Primary Analysis}

To test the primary hypotheses laid out in this PAP, I will utilize and Ordinary Least Squares (OLS) regression. In order to test H1, I will utilize the subsample of respondents that were assigned to control and those assigned to the racial diversity frame with the following OLS specification:

\begin{equation}
    Trust = \beta_0 + \beta_1 D + \beta_2 Black\_Respondent +  \beta_3 Black\_Respondent\times D  + \gamma \mathbf{X} + \epsilon
\end{equation}\ref{eq:h1}

Where $Black\_Respondent$ is self-reported respondent race, $D$ is treatment status, $X$ are covariates, and $\epsilon$ is the error term. The constant term, $\beta_0$, gives the baseline evaluation for the control group among the excluded category for measured covariates. $\beta_1$ gives the treatment effect of observing a racial diversity focused hiring drive, and $\beta_2$ is the difference in treatment effects between Black and non-Black respondents.

\subsection{Exploratory Covariate Analysis}

In order to explore descriptive patterns and interactions between covariates and treatments, I will utilize a non-parametric approach using random forest machine learning models. This approach is preferred over traditional linear specifications for several reasons. First, random forests are naturally suited for capturing high-dimensional interactions and non-linearities without requiring the researcher to pre-specify the functional form (\cite{wager2018estimation}). Second, by using a machine learning framework, I can systematically explore treatment effect heterogeneity while mitigating the multiple testing concerns inherent in exhaustive subgroup analyses (\cite{imai2013estimating}).

\subsection{Open-Ended Response Analysis}

I plan to analyze open-ended responses using text-analysis approaches to measure latent values of trust in policing. By conducting a series of pairwise comparisons between respondent open-ended responses, it is possible to construct a measure of the latent value of trust for each respondent using the Bradley-Terry pairwise-comparison model (\cite{bradley1952rank}). The Bradley-Terry model is a pairwise-comparison model where a 'winner' is selected among each pair-comparison with the assumption that, as the difference between the underling parameter of interest grows between the 'players' the probability of the winner being correctly classified also increases. More formally, the assumption is that for a pair of players $i$ and $j$, $\text{logit}[\text{pr}(i \text{ beats } j)] = \lambda_i - \lambda_j$, where $\lambda$ is the underlying latent parameter of interest. Detailed introductions to the Bradley-Terry model are available in \cite{bradley198414} and \cite{agresti2013categorical}.

In the context of trust in policing, it is possible to use the open-ended responses to construct a latent measure of respondent's underlying trust as an outcome measure by conducting pairwise comparisons across the pool of open-ended responses. Each pair is assessed for which is more 'trusting' of the police, and, under the assumption that classification is more likely to be correct as the differences in underlying trust increases, the Bradley-Terry model can be implemented to estimate a 'trust' outcome measure based on the qualitative data. This approach can be implemented using the \textbf{BradleyTerry2} package in R (\cite{turner2012bradley}).

In practice, conducting the number of pairwise-comparisons necessary to construct this latent scale is not practical for human coders, as, for a dataset with over 1,000 respondents, assuming no missingness in open-ended responses, the possible set of pairwise comparisons is $C^{1,000}_2$, which would not be a feasible number of pairwise-comparisons for a human coder. However, using Large-Language Models (LLMs) allows for an expansion of the scope of these pairwise comparisons through the use of automated coders.

\subsection{Mechanism Tests}

Using a similar approach, it is possible to construct a mechanism test with the open-ended responses. Rather than rating the responses directly for trust, the responses will be rated for plausible causal mechanisms. There are a variety of possible causal chains that could be activated by the treatments. First, it is possible that there is an identity activated effect, whereby the presence of Black or female officers drives changes in expectations. Second, it is possible that the effects are driven by reform intent. By signaling any type fo reform, the department indicates a willingness to change their practices, leading to a more trustworthy department. Finally, it is possible that cynicism drives the result. The signal of reform, be it more female officers or more Black officer, is seen as an empty promise and trust remains the same or decreases as a result of the signal of reform. Each of these mechanisms can be identified using text analysis, so the same Bradley Terry model can be applied for measurement of mechanism tests. 

\subsection{Validation}

In order to validate the work of the LLM, I plan to utilize a human-in-the-loop validation scheme. 

\subsection{Limitations and Threats to Inference}

\section{Conclusion}


\newpage
\bibliographystyle{apsr}\singlespacing
\renewcommand{\refname}{References\vspace{.3cm}}
\bibliography{Policing}
%\printbibliography




\end{document}
